\section{Related Work and Positioning}\label{sec:related}

\subsection{Time-modulated surfaces and OFDM}
Space-time-coding digital metasurfaces~\cite{zhang2018} control harmonics by the
switching pattern: a time delay of an element's pattern multiplies harmonic $h$ by a
linear phase, which steers harmonic beams. In communications, time-modulated
intelligent reflecting surfaces aligned with OFDM symbols have been designed for
waveform security, so that the intended direction receives an undistorted signal and
other directions receive scrambled symbols~\cite{xu2023tmirs}; pattern design has been
cast as generative search~\cite{tao2025gflownet,wang2026generative} and extended to
sensing~\cite{tao2025isac}. In this line of work the legitimate receiver is a
conventional OFDM receiver and the harmonic coupling it sees is either designed to be
small in its direction or treated as interference. All of it uses an ideal,
frequency-flat element. The only work that undoes the mixing is an attack: an
eavesdropper recovers directional-modulation symbols by blind independent component
analysis, resolving ambiguities with the Toeplitz structure that the ideal model
implies~\cite{tao2025secure}. Rapidly time-varying RIS have also been analysed without
OFDM~\cite{verde2023}. Iudice \emph{et al.}~\cite{iudice2026} gave the first physically
consistent model with dispersive elements, showed that dispersion breaks the Toeplitz
structure, and listed equalization across coupled subcarriers as open.

Closest in spirit outside RIS are OFDM backscatter systems, in which a tag switches its reflection within OFDM
symbols: the receiver there decodes the \emph{tag's} data from the switching, or avoids the resulting intercarrier
interference with unused subcarriers~\cite{liu2021tscatter,bsici2024,sr2024free}; space-time-coding metasurfaces
have also been used for sensing from harmonic amplitudes and for non-coherent modulation~\cite{stcisac2025,stcdiff2025}.
Linear periodically time-varying channels are well studied in power-line communication, but with periods much
longer than a symbol and hence no harmonic coupling at the subcarrier spacing~\cite{plc_lptv}.

\emph{Difference.} We design the legitimate receiver for the dispersive model: we
estimate and equalize the coupling instead of treating it as interference, and our
estimator does not rely on a Toeplitz (ideal-element) structure, which dispersion breaks.

\subsection{Estimating non-diagonal OFDM channels}
Intra-symbol channel variation from Doppler makes the OFDM channel matrix
non-diagonal. The standard estimator expands the time variation on a basis: complex
exponentials (CE-BEM)~\cite{tang2007}, or discrete prolate spheroidal
sequences~\cite{zemen2005}, giving a banded frequency-domain matrix with one unknown
gain profile per harmonic. Sparse recovery over harmonic--delay dictionaries
(orthogonal matching pursuit~\cite{tropp2007}) is an alternative. These methods assume
that the time variation is \emph{unknown}; they must estimate it.

\emph{Difference.} For a TM-RIS the time variation is known, up to the propagation
and the element response. Our estimator uses the known timing and estimates only the
static propagation (with the dispersion absorbed). The difference is large in
practice: with two pilot symbols CE-BEM reaches 2.2/1.1\,bit/s/Hz at 20/30\,dB against
our 6.5/9.7, and a full-harmonic CE-BEM needs 32 pilot symbols to approach our
two-pilot result (Section~\ref{sec:res-est}). We use CE-BEM, with the same delay prior
as ours, as the main baseline.

\subsection{Equalization of non-diagonal OFDM channels}
Banded MMSE equalizers keep only the strongest neighbouring subcarriers~\cite{schniter2004,rugini2005}.
Iterative Krylov solvers (LSQR, GMRES) with fast products through product-convolution
operators equalize a basis-expansion channel without forming the matrix~\cite{hrycak2010}.

\emph{Difference.} We adopt the iterative approach. For a TM-RIS the product-convolution
form is exact (Proposition~\ref{prop:sep}), not a basis-expansion approximation, and for
identical patterns all elements share one kernel per state, so a product needs one
batched FFT per state. Banded equalizers fail here even with the true channel, because
of the $1/h$ harmonic envelope (Section~\ref{sec:res-est}).

\subsection{RIS channel estimation}
Conventional RIS channel estimation assumes a configuration that is constant during
each OFDM symbol and changes it between symbols, so that the cascaded channel can be
separated~\cite{zheng2020ofdm}; low-rank and sparse structure of the cascaded channel is
widely exploited~\cite{wang2020compressed}. Two-timescale schemes estimate the
quasi-static BS--RIS channel rarely and the mobile RIS--user channel often~\cite{hu2019twotimescale}.

\emph{Difference.} None of these handles switching inside the symbol, where each
received subcarrier mixes many transmitted ones. We use low-rank structure and the
two-timescale principle, but inside the harmonic-coupled model, where the static side
also carries the element dispersion.

\subsection{Synchronization}
Timing offsets between RIS reflection patterns and the receiver have been studied for
distributed multi-RIS systems with maximum-likelihood estimators~\cite{zhao2022multiris}
and for RIS-assisted symbiotic radio, where the RIS modulates its own data and an
unknown offset causes intercarrier interference that is estimated blindly from virtual
subcarriers~\cite{chen2024sr}. Bayesian model selection and model averaging over timing
offsets have been used in OFDM channel estimation~\cite{access2019bayes,hku2010}, building
on the evidence framework~\cite{mackay1992}.

\emph{Difference.} We show that for a TM-RIS a clock offset is exactly an opposite
delay shift of the two hops. Consequently the pilot residual is nearly flat in the
offset, a fact that explains the failure of residual search at high SNR and motivates
evidence-based estimation. We use the evidence of our own first-stage model; the tool
is known, its necessity and its use here are not.

\begin{table}[!htbp]\centering
\caption{Capabilities of existing approaches and of the proposed receiver for OFDM through a TM-RIS.
``Pilots'' is the number of OFDM pilot symbols needed to come within 0.3\,bit/s/Hz of perfect channel knowledge at 20\,dB in our setting, pre-registered fair rerun with 30 channels and the robust equalizer for every estimator (Table~\ref{tab:fair}); $>32$: not reached with 32.}
\label{tab:features}
\resizebox{\textwidth}{!}{%
\begin{tabular}{p{4.4cm}ccccccc}\toprule
Approach & \parbox{1.3cm}{\centering in-symbol\\ switching} & \parbox{1.2cm}{\centering dispersive\\ element} & \parbox{1.4cm}{\centering element\\ model needed} & \parbox{1.3cm}{\centering full\\ coupling} & pilots & \parbox{1.3cm}{\centering clock\\ offset} & \parbox{1.6cm}{\centering cost per\\ $\bm H\bm x$} \\\midrule
One-tap OFDM (TM-IRS designs~\cite{xu2023tmirs,tao2025gflownet}) & coupling as noise & -- & -- & no & 1 & -- & $M$ \\
Blind ICA attack~\cite{tao2025secure} & yes & no (Toeplitz) & ideal & partly & 0 & -- & -- \\
Model-aware LMMSE (ours, Sec.~\ref{sec:res-model}) & yes & yes & exact & yes & 2 & -- & $QM\log M$ \\
Ideal-element LMMSE (Table~\ref{tab:fair}) & yes & no & ideal & yes & $>32$ & -- & $QM\log M$ \\
CE-BEM, $|h|\le 12$~\cite{tang2007} & generic & yes & no & truncated & $>32$ & -- & $HM$ \\
CE-BEM, all harmonics~\cite{tang2007} & generic & yes & no & yes & $>32$ & -- & $M^2$ (dense), $M\log M$ per shift (CG) \\
OMP, harmonic$\times$delay~\cite{tropp2007} & generic & yes & no & truncated & $>32$ & -- & sparse \\
RIS-OFDM estimation~\cite{zheng2020ofdm} & no & -- & -- & -- & -- & -- & $M$ \\
Two-timescale RIS~\cite{hu2019twotimescale} & no & -- & -- & -- & -- & -- & -- \\
\textbf{Proposed (timing-only)} & \textbf{yes} & \textbf{yes} & \textbf{no} & \textbf{yes} & \textbf{2 (TDL-C: 4)} & \textbf{blind} & $\bm{QM\log M}$ \\\bottomrule
\end{tabular}}
\end{table}
